\originalsection{从推迟核到移动点源的雅可比}
先把四维推迟积分中的时间变量保留, 点源的空间积分完成后, 标势成为
$V=q(4\pi\eps)^{-1}\int\delta[t-t'-R(t')/c]\dd t'/R(t')$.
对 $h(t')=t-t'-R(t')/c$, 有 $\dd R/\dd t'=-\nn\cdot\vct v$, 所以在根 $t_r$ 处 $h'=-1+\nn\cdot\vct\beta=-\kappa$. 亚光速保证 $\kappa>0$, 因而简单根分布公式给 $\delta(h)=\delta(t'-t_r)/\kappa$. 这产生 $V=q/(4\pi\eps\kappa R)$; 矢量势以相同雅可比给 $\vct A=\vct\beta V/c$. 观察点固定时 $h$ 对时间微分, 或观察时间固定时对位置微分, 分别得到前文的 $\partial_tt_r=1/\kappa$ 与 $\nabla t_r=-\nn/(c\kappa)$.

\originalsection{实际微分 Li\'enard--Wiechert 势}
为让速度和加速度项的分母能够逐步核对, 记 $F=\kappa R=R-\vct\beta\cdot\vct R$, $K_q=q/(4\pi\eps)$. 下文 $\dot{\vct\beta}$ 指对源时间 $t_r$ 求导, 所有量在 $t_r$ 取值. 推迟位置的导数为
\[
 \partial_t\vct R=-\frac{c\vct\beta}{\kappa},\qquad
 \partial_j R_i=\delta_{ij}+\frac{\beta_in_j}{\kappa}.
\]
对 $F=R-\vct\beta\cdot\vct R$ 用乘积法则, 注意 $\partial_t\vct\beta=\dot{\vct\beta}/\kappa$、$\partial_j\beta_i=-\dot\beta_in_j/(c\kappa)$, 得
\begin{align*}
 \partial_tF&=\frac{c(\beta^2-\nn\cdot\vct\beta)-R\nn\cdot\dot{\vct\beta}}{\kappa},\\
 \nabla F&=-\vct\beta+\frac{1-\beta^2+(R/c)\nn\cdot\dot{\vct\beta}}{\kappa}\nn.
\end{align*}
这两行同时包含 $R$ 改变、源速度改变和推迟时间改变, 任何一项都不能当成常量删掉. 由 $V=K_q/F$、$\vct A=K_q\vct\beta/(cF)$,
\[
 \ee=K_q\left[\frac{\nabla F+\vct\beta\partial_tF/c}{F^2}
 -\frac{\dot{\vct\beta}}{c\kappa F}\right].
\]
通分后, 分子中的速度部分是 $(1-\beta^2)(\nn-\vct\beta)$, 加速度部分是
$(R/c)[(\nn-\vct\beta)(\nn\cdot\dot{\vct\beta})-\kappa\dot{\vct\beta}]$,
共同分母为 $\kappa^3R^2$. 三重积恒等式将方括号写成 $\nn\times[(\nn-\vct\beta)\times\dot{\vct\beta}]$, 恰得前文的完整场式.

磁场同样直接取旋度. $\nabla\times\vct\beta=\nabla t_r\times\dot{\vct\beta}=-\nn\times\dot{\vct\beta}/(c\kappa)$, 故
$\bb=-K_q\nabla F\times\vct\beta/(cF^2)-K_q\nn\times\dot{\vct\beta}/(c^2\kappa F)$.
若把 $\nabla F=-\vct\beta+\lambda\nn$ 中的系数记 $\lambda$, 上两导数满足 $\lambda=1-\partial_tF/c$. 代回即给 $\bb=\nn\times\ee/c$. 所以这条关系对该点电荷的全场成立, 并非仅从平面波类推.

对静止电荷 $\vct\beta=\dot{\vct\beta}=0$, 速度场恢复库仑场. 对匀速电荷只有 $1/R^2$ 项; 对非相对论加速电荷, $\kappa\simeq1$、$\dot{\vct\beta}=\vct a/c$, 加速度场为 $q\nn\times(\nn\times\vct a)/(4\pi\eps c^2R)$. 它的场幅随 $R^{-1}$ 衰减, 乘球面积后仍有有限功率, 是辐射的几何原因. 本段的中文求导补足原第15讲的数学内容, 不需要读者回英文原件补步骤.
